% EGA I, Chapitre 0, §3.6 and complete3.6.1; French1794-1806.
% GPT-6.1 Sol / Ultra; source-linked draft, R601.
\subsection{Haces simples y haces localmente simples}

\begin{env}[3.6.1]
\label{0.3.6.1-fr}
Diremos que un \emph{prehaz} $\mathcal{F}$ sobre $X$, con valores en
$\boldsymbol{K}$, es \emph{constante} si los morfismos canónicos
$\mathcal{F}(X)\to\mathcal{F}(U)$ son \emph{isomorfismos} para todo
abierto no vacío $U\subset X$~; se observará que $\mathcal{F}$ no es
necesariamente un haz. Se dice que un \emph{haz} es \emph{simple}
si está asociado (3.5.6) a un prehaz constante. Se dice que un haz
$\mathcal{F}$ es \emph{localmente simple} si todo $x\in X$ admite una vecindad
abierto $U$ tal que $\mathcal{F}|U$ sea simple.
\end{env}
